{\color{ink}\bfseries 2.10 \(\cdot \) 原题前驱图；图下文字标明题干给出}\par\vspace{1bp}只需为跨线程的A\(\to \)C、C\(\to \)E各设一个初值0的信号量aDone、cDone：\par T1：A;\allowbreak{} V(aDone);\allowbreak{} P(cDone);\allowbreak{} E;\allowbreak{} F;\allowbreak{}\par T2：B;\allowbreak{} P(aDone);\allowbreak{} C;\allowbreak{} V(cDone);\allowbreak{} D;\allowbreak{}\par aDone只在A完成后发布，cDone只在C完成后发布；B\(\to \)C、C\(\to \)D、E\(\to \)F由各线程自身顺序保证，D与E可并发。两个方向不能直接共用同一个普通信号量，否则T1可能在A之后取回自己刚发布的许可，使E抢在C之前执行。
